\scanpage{430}
\mspage{169}
% archive sheet number: 429

\diagram{\begin{tikzcd}[ampersand replacement=\&,column sep=3mm,row sep=8mm]
a'\arrow[r,"u"]\arrow[d,"u"'] \& a\arrow[d,dash,"\mathrm{id}"] \&
a\arrow[l,"\mathrm{id}"']\arrow[d,"\mathrm{id}"]\arrow[r]\arrow[rrrr,dash] \&
{} \& {}\arrow[d,"\mathrm{id}"',shorten <=4pt,shorten >=4pt] \& {} \&
b\arrow[l]\arrow[r,"\mathrm{id}"]\arrow[d,"\mathrm{id}"] \& b\arrow[d,"\mathrm{id}"'] \&
b'\arrow[l,"v"']\arrow[d,"v"] \& \widetilde c\quad\text{\unsure{type} }t_{n+1}\\
\unsure{a|a'}\arrow[r,"\mathrm{id}"] \& a \& a\arrow[l,"\mathrm{id}"']\arrow[r]\arrow[rrrr,dash,"c"] \& {} \& {} \& {} \&
b\arrow[l]\arrow[r,"\mathrm{id}"] \& b\arrow[r,dash,"\mathrm{id}"] \& b \& \varphi_{n,n+1}^{*}(c).
\end{tikzcd}}

On doit seulement vérifier que l'endofoncteur $\psi\varphi$ de $C_{\mathbb N^{\unsure{\circ|o}}}(c,u,v)$ est homotope à l'identité. Mais \unsure{ce} foncteur transforme $(c',w')$ en $(\widetilde{c'}=a'_{\Delta_+\Delta_-}\,c'\,b'_{\Delta_+\Delta_{\unsure{+|-}}},\overline{w'})$, où $\overline{w'}$ est \unsure{encore|le lien} le composé

\diagram{\begin{tikzcd}[ampersand replacement=\&,column sep=9mm,execute at end picture={
\draw[-{Straight Barb[length=2pt]},rounded corners=2pt] (lc.south) -- ++(0,-3mm) -| (rc.south) node[pos=0.25,below] {$\overline{w'}$};
}]
|[alias=lc]|\widetilde{c'}\arrow[r,"\mathrm{can}"] \& c'\arrow[r,dash,"w'"] \& |[alias=rc]|c
\end{tikzcd}}

On a donc

\[
\psi\varphi\longrightarrow\mathrm{id},
\]

$\psi\varphi$ est homotope au foncteur identique.

Cela prouve la lissité de $\unsure{\mathcal C|\mathbb C}_{\mathbb N^\circ}\unsure{\mathrel{-}|\longrightarrow}X\times X$. Pour la propreté, il faut prouver la lissité dans le cas de $X^\circ$, et du système projectif $\uline{t}^{\mathrm{opp}}\simeq(t_n^{\mathrm{opp}})$. Mais on a $\uline{t}^{\mathrm{opp}}=\uline{t}$, et \deleted{\unsure{il}} suffit d'appliquer ce qu'on a déjà prouvé \deleted{\unsure{pour}} au cas de $X^\circ$, pour voir que $\mathcal C_{\mathbb N^{\unsure{\circ}}}\longrightarrow X\times X$ est propre.

Il ne reste plus que le lemme à prouver. Ça doit être \unsure{essentiellement}
